\section{Limitations and Future Work}

The current draft has a deliberately narrow soundness boundary.
The Lean development mechanizes a small core of capability, certificate, resource, and effect properties, but not yet a full source-to-target elaboration theorem.
The Python checker validates generated plans against the current sidecar format, but the implementation does not yet extract a complete machine-checkable typing derivation for every accepted plan.
Closing this gap is the most important next step for a POPL submission.

Several rules in the current case studies are marked \Assumed{}.
These rules are useful because they let the compiler exercise end-to-end acquisition strategies that connect otherwise paper-backed pieces, but they are not physics claims.
The paper therefore reports exploratory results with warnings and treats strict mode as the proof-facing policy.
Before submission, prototype bridge rules should either be replaced by literature-backed certificates or moved into a clearly separated exploratory appendix.

The layout model is also conservative.
The geometry checker verifies final rectangular patch embeddings and workspace reservations on a fixed backend grid, but it does not model complete patch routing, lattice-surgery boundary compatibility, or factory scheduling optimality.
This limitation is acceptable for the current language claim because topology and workspace violations are already visible to the checker.
It is not sufficient for claiming a production-quality FTQC layout compiler.

The benchmark set is small.
The case studies validate the calculus and checker interface, but they do not compare against production FTQC compilers or a broad suite of fault-tolerant workloads.
Future work should add larger logical programs, more backend profiles, stronger acquisition baselines, and ablations that isolate the effect of hybrid-region planning, certificate policies, and resource-factory choices.
